\documentclass[../Main.tex]{subfiles}

\begin{document}
\section{Introduction and Examples}
\begin{definition}{Transformation}
    Given a unitary operator $U$, define a \underline{transformation} of a quantum system to be either a map of states:
    \begin{equation*}
        \state \mapsto U\state,\qquad \ket{\phi} \mapsto \ket{\phi}U^\dagger
    \end{equation*}
    or a map of operators:
    \begin{equation*}
        A \mapsto U^\dagger A U
    \end{equation*}
\end{definition}
\begin{remarks}
    \item Mapping state and mapping operators are equivalent because inner products of states are unchanged. However, matrix elements do change in general:
        \begin{equation*}
            \opbraket \phi A \psi \mapsto \opbraket\phi{U^dagger A U}\psi.
        \end{equation*}
    \item Unlike changing pictures we change states or operators but not both.
\end{remarks}
\begin{definition}{Symmetry}
    A transformation of a quantum system is a \underline{symmetry} of that quantum system if:
    \begin{equation*}
        U^\dagger H U = H \text{ or } \cmmt{U}{H} = 0.
    \end{equation*}
\end{definition}
It follows that:
\begin{itemize}
    \item in the \schr picture, if $\ket{\psi(t)}$ is a solution of the \schr equation, then so is $U\ket{\psi(t)}$;
    \item in the Heisenberg picture, if $A(t)$ is a solution of the Heisenberg equation, then so is $U^\dagger A(t) U$.
\end{itemize}
Consider a group $G$ of transformations of a quantum mechanical system. Then for each $g \in G$, $U(g)$ will be a unitary operator that will act in a similar way to $g$. Indeed,
\begin{align*}
    U(g_1) U(g_2) &= U(g_1 g_2) \\
    G(e_G) &= 1 \\
    U(g^{-1}) &= U(g)^{-1} = U(g)^\dagger
\end{align*}
where $e_G$ is the identity element of the group $G$. Our aim is to find $U(g)$ when $G$ is a group of translations, rotations or reflections, because in these cases we know how $G$ should act geometrically.

From this, we can infer the action of $U(g)$ on position eigenstates, $U(g) \ket{\vec x} = \ket{g(\vec{x})}$, or the position operator:
\begin{equation*}
    U(g)^\dagger \opvec x U(g) = g(\opvec{x})
\end{equation*}
These states are equivalent since:
\begin{equation*}
    \opvec{x}(U(g)\ket x) = U(g) g(\opvec{x}) \ket{\vec x} = g(\vec x) U(g) \ket{\vec x}
\end{equation*}
On a general state,
\begin{equation*}
    \state \mapsto U(g) \state
\end{equation*}
where the position space wavefunction of our new state is:
\begin{align*}
    \opbraket{\vec x}{U(g)}{\psi} &= \left(U(g)^\dagger \ket{\vec x}\right) \state \\
    &= \left(U(g^{-1}) \ket{\vec x}\right) \state \\
    &= \inn{g^{-1}(\vec{x})}{\psi}
\end{align*}
\begin{example}[Translations in one dimension]
    We have $G = \{g_a\}_{a \in \R}$ and:
    \begin{equation*}
        x\mapsto g_a(x) = x + a
    \end{equation*}
    representing a translation. Write $U(a)$ to mean $U(g_a)$. Then:
    \begin{equation*}
        U(a) \ket x \mapsto \ket{x + a}
    \end{equation*}
    or:
    \begin{equation*}
        U(a)^\dagger \hat x U(a) = \hat x + a
    \end{equation*}
    and we know that the operator with these properties is:
    \begin{equation*}
        U(a) = \exp\left(-\frac{ia}{\hbar} \hat p\right)
    \end{equation*}

    The effect of translation by $a$ on the wavefunction $\psi(x)$ is to generate the new wavefunction $\psi(x - a) = \psi(g^{-1}_a(x))$.
\end{example}
\end{document}